\documentclass[conference]{IEEEtran}
\usepackage{cite}
\usepackage{url}
\usepackage{hyperref}
\usepackage{xspace}
\newcommand{\privoke}{\textsc{PriVoke}\xspace}
\IEEEoverridecommandlockouts
\begin{document}
\title{\privoke: Interim Evidence for a Client-Side Prompt Privacy Firewall}
\author{
\IEEEauthorblockN{Nicholas Bisset\IEEEauthorrefmark{1}}
\IEEEauthorblockA{Monash University Clayton, Australia\\nbis0006@student.monash.edu}
\and
\IEEEauthorblockN{Manpreet Sahadevan, Nixon Ngai, and Sanoop Mallissery}
\IEEEauthorblockA{Monash University Malaysia, Malaysia\\
\{msah0018, nnga0004\}@student.monash.edu, sanoop.mallissery@monash.edu}
\thanks{Internal research progress draft: evidence checkpoint 5 October 2026.}
\thanks{\IEEEauthorrefmark{1}Corresponding author: Nicholas Bisset.}
\thanks{\href{https://github.com/Densetsu152637/Privoke-Research-Project}{Research repository}.}
}
\maketitle
\begin{abstract}
Prompts submitted to hosted large language models can disclose sensitive information before provider-side controls intervene. \privoke combines browser request interception with rules, named entity recognition and a local streamed classifier. This interim report evaluates layer contributions, bounded classifier-head updates and a separate annotation-presence gate. On 502 development sentences, the initial full pipeline achieved 93.56\% annotation-presence recall and 19.75\% specificity. A later negative-curriculum checkpoint achieved 90.53\% recall and 29.41\% specificity. An opt-in efficient presence gate on the revised original pipeline reduced false positives from 185 to 63, while adding five misses, giving 91.67\% recall and 73.53\% specificity. Initial template updates did not improve pipeline recall; none of nine sparse-head updates qualified for retention. These findings expose tradeoffs rather than establish contextual privacy or safe actions. The 498-row final partition remains locked and unscored; independent contextual adjudication and complete installed-extension measurements remain unresolved.
\end{abstract}
\begin{IEEEkeywords}
Prompt privacy, client-side enforcement, annotation presence, controlled model updates, local differential privacy.
\end{IEEEkeywords}

\section{Introduction}
\label{sec:introduction}
Hosted large language models (LLMs) support writing, coding and analysis, but prompt submission moves user-supplied content across an external disclosure boundary. Privacy-sensitive applications require attention to what is submitted, as well as what a model returns~\cite{priyanshu2023chatbots,barbera2025ai}. Workplace and public-sector incidents illustrate accidental disclosure of confidential code and sensitive case information~\cite{SamsungBans2023,ChatGPTChildProtection}. Provider retention controls do not replace an endpoint decision about whether to transmit a particular request.

Client-side prompt sanitization has close precedents. Casper combines regular expressions, NER and a local model in a browser privacy pipeline~\cite{Casper2024}; HaS and ProSan investigate other prompt-protection approaches~\cite{HaS2024,ProSan2024}. We credit this architectural overlap. Our question is whether bounded updates and a separately learned annotation-presence control improve the detection tradeoff in the implemented \privoke pipeline. Combining familiar components alone does not establish novelty or superiority.

\privoke warns or cancels supported requests; WARN forwards the original request and BLOCK cancels it. It does not rewrite prompts. Its streamed classifier uses repository-owned, randomly initialized compact encoder weights and fitted heads, rather than a pretrained conversational LLM. Locally randomized telemetry is a separate monitoring mechanism, not a training source.

This progress report contributes (1) an implementation-grounded account of interception, detection, model publication and telemetry boundaries; (2) matched development measurements of layer ablations and controlled head-update regressions; and (3) measurements of a separate sparse presence model and exploratory semantic gate, retaining action failures and data limitations. These are interim observations, not submission-readiness or universal privacy claims.

\section{Background and Related Work}
\label{sec:background}
Rules identify structured fields, while NER identifies entity types in unstructured text. Entity presence and contextual privacy differ: a public name can be annotated as PII, while private business information may lack an annotated entity. Prompt sanitization can redact or transform content~\cite{Casper2024,HaS2024,ProSan2024}; \privoke applies ALLOW, WARN and BLOCK decisions. No matching Casper experiment or user study was performed, so we make no numerical ranking or usability claim.

Memorization studies show that training information can sometimes be reconstructed~\cite{haim2022reconstructing}; they do not show that every submitted prompt is used for training. Controlled template testing can diagnose failures, but accepted parameter updates and internal guards do not establish generalization.

Local differential privacy (LDP) randomizes reports before collection~\cite{cormode2018privacy}. For neighboring event values $x,x'$ and output $o$, pure $\varepsilon$-LDP requires $\Pr[M(x)=o]\leq e^\varepsilon\Pr[M(x')=o]$. Event-value protection does not conceal report presence or transport metadata.

\section{Threat Model and Scope}
\label{sec:scope}
The protected asset is prompt text on supported configured browser endpoints. We consider accidental disclosure and bounded textual evasion, trusting the browser, extension, workstation supervisor, detector, artifacts and configured service boundary. A compromised page can interfere with same-window decision messaging; this is not an isolated defense against a compromised client. Unsupported formats, nontext attachments, alternative transports, provider-side retention and bypass outside intercepted paths are outside the demonstrated scope. Optional hosted inference sends text to another service and has a different disclosure boundary.

The hook supports bounded POST fetch and asynchronous XMLHttpRequest paths. Requests without extracted prompt text can pass without analysis. Bridge and transport errors block, but known semantic unavailability omits that layer if regex or NER remains enabled; semantic-only use blocks. A requested-layer error upgrades an otherwise ALLOW result to BLOCK while retaining surviving WARN or BLOCK decisions.

\section{Design and Implementation}
\label{sec:system_architecture}
\subsection{Runtime and Classification}
The extension uses a workstation supervisor with loopback lifecycle control on port 50056, detector on 50057 and gRPC-Web bridge on 8080. The Compose detector on internal port 50054 is an independent process used by the fuzzer and evaluators, not an extension fallback.

Normalization applies NFKC, lowercase conversion and bounded deobfuscation while preserving provenance to original-text spans. Regex runs first by default and short-circuits remaining requested layers on BLOCK. Otherwise pending layers run concurrently. The spaCy \texttt{en\_core\_web\_sm} NER path maps PERSON, GPE, LOC, FAC and ORG, not DATE. Fusion retains strongest sensitivity, strongest known visibility and category union. Successful evidence survives other layer errors, and fusion does not weaken individual actions.

Classification uses sensitivity S0--S3, visibility P0--P4 and PU, and flags for health, politics, religion, criminal, financial, sexual, child, location, identity and third-party information. S3 blocks and S2 warns; identity/location with restricted visibility or both categories can warn. Confidence below 0.5 moderates S3 to WARN and non-identifying S2 to ALLOW. Binary annotated-entity prediction is distinct from action.

\subsection{Streamed Contextual Classifier}
The encoder hashes tokens into a bounded vocabulary and truncates context. Seeded random initialization supplies embeddings, attention and feed-forward weights; 43 authored bootstrap examples fit only six sensitivity, visibility and category head tensors. Runtime updates also alter heads only. Efficient, balanced and quality configurations use respectively one/two/three layers, hidden dimensions 24/32/32, context limits 64/96/128 and vocabularies 512/512/768. Their names do not establish quality.

The Go service returns one immutable artifact snapshot per request, not a continuing subscription. Clients validate chunk order, offsets, shapes, identity and completeness, reconstruct models and retain a short-lived cache. Local inference supports CPU with CUDA/MPS paths where available. The \texttt{latest} alias selects a configured current artifact, not a historical version.

\subsection{Training and Publication}
The fuzzer deterministically samples labeled templates or a configured curriculum and can apply fixed transformations; it has no generative sampling-temperature control. Automatic/background cycles are disabled by default; studies explicitly request cycles. It sends examples to the runtime to compute bounded head deltas. It does not download tensors or train the encoder. Unlabeled examples may use the model's predictions as targets and are not independent supervised truth.

The updater validates deltas, rejects stale base versions and serializes accepted updates. Atomic replacement includes the committed receipt; durable historical receipts support replay recovery without applying an identified payload twice. Candidate checks use publication-compatible float32/clipping arithmetic and an internal held-out guard. This is sequential centralized updating, not federated learning. Compose internal gRPC is plaintext; deployment outside a trusted service network requires the configured authenticated transport boundary.

\subsection{Separate Presence Task and Training Mechanics}
A typed annotation-presence RPC exposes a separate word/character TF-IDF logistic model. Efficient/balanced/quality maximum capacities are 2,000+2,000, 8,000+8,000 and 16,000+16,000 features. Vocabulary and IDF remain frozen during head updates; coefficient blocks contain at most 4,096 values. Presence supplies no sensitivity, visibility, category or action truth.

The ordinary contextual path remains unchanged. An explicit research gate on the balanced contextual model suppresses only semantic findings when presence is ABSENT, preserving regex and NER. ABSENT does not mean safe. Callers must verify the returned trace rather than assume an older server applied the option.

Separate offline CPU trainers implement head-only and full-encoder supervised mechanics for contextual and scratch binary transformers: explicit targets, finite-gradient checks, clipping, rollback and export. Presence inference also supports six explicit scratch model IDs. These mechanics differ from measured sparse gradient dispatch. No research-data scratch-transformer fit or quality result is reported.

\subsection{Telemetry}
Workstation telemetry is opt-in through runtime configuration; Compose evaluation settings are separate. Five fields are randomized: action, risk bucket, primary category, normalized model release and coarse UTC time-of-day. Generalized randomized response uses default event budget $\varepsilon=1$, equally divided across fields, with $\delta=0$. A durable installation-local ledger limits cumulative expenditure to $\varepsilon=8$ per UTC day and suppresses emission when exhausted or unusable.

This assumes a trusted emitter randomizes values and preserves its ledger. The collector validates schema/budget but cannot prove randomization. Reports contain no raw prompt, direct user identifier, exact timestamp or exact risk score. Marginal estimates apply inverse randomized-response correction and clipping, which can introduce bias. Exact sample counts, report presence and network timing/metadata remain outside the guarantee. Reports do not train models; telemetry utility is unmeasured.

\section{Methods}
\label{sec:methods}
We ask RQ1: how do layers contribute to recall and false alarms? RQ2: do bounded updates improve both class rates? RQ3: how do presence inference and gating trade errors and actions? RQ4: what detector costs are measured?

English PIIMB data\footnote{\url{https://huggingface.co/datasets/piimb/pii-masking-benchmark}} are pinned to revision \texttt{4a13e9ffe6fd0d275efbde8afd4d8d8f1ffc2133}. Seed 3102026, deduplication, conflict exclusion and source-document grouping select 1,000 rows: 502 development (264 annotated-positive, 238 annotation-negative) and 498 final (236/262). Final remains locked and unscored. Later controls fit 3,832 training rows and calibrate on 968 validation rows (475/493), retaining the same development population. Fitting on unused documents from this upstream benchmark-test corpus creates a custom within-corpus study, not an untouched official test score.

A prediction is positive if sensitivity is S1--S3 or any category is set, regardless of action. Annotations do not establish exact spans, contextual privacy, visibility or safe actions. Provided labels are retained even when unannotated text triggers broader policy rules. Precision in this class mix is not deployment precision.

Comparisons match IDs, labels and groups; reported 502-row detector runs have zero runtime errors. Independent Presidio\footnote{\url{https://microsoft.github.io/presidio/}} uses English \texttt{en\_core\_web\_sm} and threshold 0.5, smaller than its upstream default NLP model. Initial updates use three independent seeds, 256 source prompts and rate 0.03. Later negative-curriculum selection maximizes specificity subject to at least 90\% recall and recorded floor/tie rules. Validation selects sparse C values, thresholds and gates before development scoring. Earlier development findings informed these designs.

Computed paired source-group bootstrap intervals use 2,000 resamples, seed 3102026 and 465 development groups. They are descriptive after search/selection, not confirmatory significance. Cascade intervals remain uncomputed. The practical 90\% recall/90\% specificity targets are development criteria, not publication requirements. The companion ledger separately records artifact checksums, file hashes, parameter fingerprints, source/image identities and raw reports.

\section{Measured Development Results}
\label{sec:results}
\subsection{RQ1: Layer Contributions}
Immutable original v0.3.0 comparisons at source \texttt{712ed72} show increased pipeline recall with many annotation-negative detections (Table~\ref{tab:ablations}). This cannot identify contextual false alarms without independent labels.
\begin{table}[t]
\centering
\caption{Initial matched PIIMB development ablations, $n=502$ (264 positive, 238 negative). Counts are TP/TN/FP/FN; rates are percentages.}
\label{tab:ablations}
\begin{tabular}{lrrr}\hline
Configuration & Counts & Recall & Specificity\\\hline
Regex & 179/208/30/85 & 67.80 & 87.39\\
NER & 55/232/6/209 & 20.83 & 97.48\\
Regex + NER & 198/203/35/66 & 75.00 & 85.29\\
Semantic & 169/60/178/95 & 64.02 & 25.21\\
Full pipeline & 247/47/191/17 & 93.56 & 19.75\\
Presidio (small) & 208/189/49/56 & 78.79 & 79.41\\\hline
\end{tabular}
\end{table}
Later rules produce a distinct original-model baseline, 247/53/185/17 (93.56\% recall, 22.27\% specificity). It must not replace the initial ablation baseline.

\subsection{RQ2: Update Tradeoffs}
Initial transformed and ordinary rate-0.03 updates give no pipeline recall gain across three seeds and reduce specificity by 0.84--1.26 percentage points. Accepted internal checks do not imply public improvement. A historical lower-rate/rule checkpoint corrected seven initial false positives without changing recall, still giving only 22.69\% specificity.

A broader curriculum of 2,400 public negatives plus bootstrap anchors yields Table~\ref{tab:updates}. Nine independent attempts produce six scored candidates and three held-out regression rejections. Larger accepted rates fall below the recall floor. Cycle two of the selected seed reaches 236/78/160/28 (89.39\% recall); the protocol stops and restores cycle one.
\begin{table}[t]\centering
\caption{Revised-rule original and selected negative-curriculum development checkpoints, same 502 rows; rates are percentages.}
\label{tab:updates}
\begin{tabular}{lrrr}\hline
Checkpoint & TP/TN/FP/FN & Recall & Specificity\\\hline
Original & 247/53/185/17 & 93.56 & 22.27\\
Selected & 239/70/168/25 & 90.53 & 29.41\\\hline
\end{tabular}
\end{table}
Relative to the same-source original, selection removes 17 false positives and adds eight false negatives. Recall changes by $-3.03$ percentage points (descriptive 95\% interval $[-5.26,-1.12]$), specificity by $+7.14$ ($[4.00,10.48]$). Semantic recall falls from 64.02\% to 53.03\%; other layers retain some detections. This is a tradeoff, not improved semantic understanding.

\subsection{RQ3: Presence and Gating}
The separate sparse task gives Table~\ref{tab:presence}. An offline larger text control yields 248/186/52/16 (93.94\% recall, 78.15\% specificity), not a deployed contextual result. None of nine accepted sparse fuzzer updates is retained because validation specificity does not strictly improve.
\begin{table}[t]\centering
\caption{Separate annotation-presence RPC results on 502 development rows, not full-pipeline or action scores; rates are percentages.}
\label{tab:presence}
\begin{tabular}{lrrr}\hline
Profile & TP/TN/FP/FN & Recall & Specificity\\\hline
Efficient & 249/169/69/15 & 94.32 & 71.01\\
Balanced & 247/181/57/17 & 93.56 & 76.05\\
Quality & 247/189/49/17 & 93.56 & 79.41\\\hline
\end{tabular}
\end{table}
An expanded 19,993-row fit selects C=10 on validation. Each profile detects 1,000/1,000 Nemotron\footnote{\url{https://huggingface.co/datasets/nvidia/Nemotron-PII/tree/b70ffaf5ff39e079776134c5bf4381f00a9fd1ed}} and 999/999 Meddies\footnote{\url{https://huggingface.co/datasets/Meddies/meddies-pii/tree/6a5c8f5441e3b421d983c9741770262365acdd77}} source-heldout positives. Positive-only sets cannot estimate specificity. Mixed validation specificity is 66.94\%, 80.32\% and 77.89\%; no contextual inference follows.

Original-balanced gating reduces false positives but adds five misses per profile (Table~\ref{tab:cascade}). Thresholds are validation-selected and paired intervals uncomputed. Trained-control arms are skipped because validation recall is 426/475 (89.68\%), below the floor. No gate is promoted.
\begin{table}[t]\centering
\caption{Original balanced pipeline with validation-selected semantic gates, same 502 development rows; regex/NER remain intact; rates are percentages.}
\label{tab:cascade}
\begin{tabular}{lrrr}\hline
Gate & TP/TN/FP/FN & Recall & Specificity\\\hline
None & 247/53/185/17 & 93.56 & 22.27\\
Efficient & 242/175/63/22 & 91.67 & 73.53\\
Balanced & 242/169/69/22 & 91.67 & 71.01\\
Quality & 242/159/79/22 & 91.67 & 66.81\\\hline
\end{tabular}
\end{table}
The returned regex/NER union has 13 development false positives, an observed specificity ceiling of 94.54\%. This necessary bound does not show a gate can reach it while retaining recall or safe actions.

A separate 48-case authored fixture has 24 controls, 17 disclosure candidates and seven ambiguous exclusions. Labels are assistant-provisional; professor confirmation is pending. Ordinary actions are rubric-compatible on 30/41 evaluable cases, versus 32/41 for efficient and 31/41 for balanced/quality gates. All meet minimum private action on 15/17. Two required-BLOCK cases remain below BLOCK, one ALLOW and one WARN. This is not policy-safety certification.

\subsection{RQ4: Detector Costs}
Table~\ref{tab:timing} reports sequential CPU measurements on an x86\_64 Intel Core i7-13700KF with 24 visible CPUs. Contextual timings use returned evaluator elapsed time; sparse timings are internal service measurements. Summaries include the first request and exclude browser/native-messaging/bridge overhead. They are not robust cold-start or deployment-latency measurements.
\begin{table*}[t]\centering
\caption{Detector median/p95 latency (ms), 502 requests per configuration. Different timing boundaries preclude additive comparison.}
\label{tab:timing}
\begin{tabular}{lrrr}\hline
Profile & Semantic & Pipeline & Sparse internal\\\hline
Efficient & 6.45/7.58 & 229.48/289.24 & 2.524/3.932\\
Balanced & 13.45/16.91 & 267.21/338.31 & 9.921/13.803\\
Quality & 20.05/23.96 & 276.94/354.27 & 15.396/19.086\\\hline
\end{tabular}
\end{table*}
Profiles differ in initialization, depth, context, vocabulary and bootstrap epochs. Quality has higher semantic recall but lower pipeline recall/specificity than balanced. No causal capacity effect, memory result or usability benefit is established.

\section{Discussion and Limitations}
\label{sec:limitations}
Layer union raises recall while retaining false positives. Negative-head updates shift the tradeoff; presence gates suppress semantic findings, including annotated positives. The evidence does not support monotonic learning, a safe ABSENT interpretation or the joint 90/90 target. Representation/task mismatch and source-label differences remain competing explanations.

Construct validity is limited by annotation presence as a disclosure proxy and by provisional fixture labels. Internal validity is limited by development-informed changes and differing configurations. The selected trained contextual curriculum overlaps later validation, so its diagnostic comparisons are not independent. External validity is limited to English corpus samples, source imbalance and authored fixtures. Intervals after selection are descriptive.

Detector scores do not establish transmitted-request outcomes. Installed-extension validation was attempted but did not complete successfully; no complete end-to-end latency/resource evidence is available at this checkpoint. Partial forwarding probes cannot certify broad coverage. Unsupported requests, failed extraction, WARN forwarding and trusted-page assumptions bound enforcement. Event LDP excludes model training and network metadata; telemetry utility is unmeasured.

\section{Ethics, Artifacts and Remaining Work}
\label{sec:artifacts}
PIIMB documentation lists CC BY-NC 4.0 and upstream sources; raw redistribution requires license review and attribution. Saved local evidence contains prompt text and needs appropriate access controls. The companion \path{paper/research/manuscript-progress-20261005.md} maps tables to source/raw evidence and distinguishes file hashes from model identities. Ignored local artifacts are not automatically a public reproducibility package.

Generative AI assisted this revision, evidence organization and provisional fixture construction. The assistance record identifies configured model and scope; authors must verify factual content, labels, references and venue-specific disclosure. No independent human agreement or ethics approval is implied.

Final scoring must follow frozen model/rule/protocol decisions and preserve the 498-row lock until then. Independent contextual/action adjudication, complete installed-extension enforcement/cost assessment, representative deployment evidence and telemetry utility remain pending. A clean-data augmentation structural scan covers 104,728 rows, but reviewed broad-target eligibility, fresh quota-feasible partitions, fitting and scoring remain pending. Offline encoder-training mechanics need a separate research-data fit before quality claims.

\section{Conclusion}
\label{sec:conclusion}
\privoke implements interception, layered analysis, local streamed inference, governed head updates and opt-in event-level LDP telemetry. Development measurements show substantial false alarms and recall--specificity tradeoffs. Initial updates do not improve pipeline recall; later negative training and semantic gating reduce false positives with lost detections. Presence-only scores and provisional fixtures do not establish contextual safety. Final remains unscored, and broader deployment/privacy conclusions require additional evidence.
\bibliographystyle{ieeetr}
\bibliography{ref}
\end{document}
